\documentclass[aspectratio=169, 11pt]{beamer}
\usepackage[utf8]{inputenc}
\usepackage[T1]{fontenc}
\usepackage[french]{babel}
\usepackage{graphicx}
\usepackage{tikz}
\usepackage{xcolor}
\usepackage{hyperref}

\usetheme{Madrid}
\usecolortheme{beaver}

\setbeamercolor{frametitle}{fg=black,bg=white}
\setbeamercolor{title in head/foot}{fg=black,bg=white}

\setbeamertemplate{footline}{%
  \leavevmode%
  \hbox{%
    \begin{beamercolorbox}[wd=.82\paperwidth,ht=2.25ex,dp=1ex,leftskip=1.2ex]{author in head/foot}%
      \usebeamerfont{author in head/foot}%
      \footnotesize RVP2 --- Moteur de recherche multimodal (Ping) --- Barbier / Ledoux / Pascal / Wone
    \end{beamercolorbox}%
    \begin{beamercolorbox}[wd=.18\paperwidth,ht=2.25ex,dp=1ex,right,rightskip=1.2ex]{page number in head/foot}%
      \usebeamerfont{page number in head/foot}%
      \footnotesize \insertframenumber/\inserttotalframenumber
    \end{beamercolorbox}%
  }%
}

\title[Recherche semantique]{Fonctionnement de la recherche semantique}
\author{Nathan Barbier - Projet Informatique}
\date{\today}

\begin{document}

\begin{frame}{Recherche semantique : construction des embeddings}
    {\small
    \begin{block}{Pipeline offline}
        \begin{itemize}
            \item CSV point par point
            \item description textuelle enrichie
            \item encodage par \texttt{paraphrase-multilingual-MiniLM-L12-v2}
            \item vecteur dense de dimension \textbf{384}
        \end{itemize}
    \end{block}

    \begin{block}{Contenu de la description}
        \begin{itemize}
            \item joueurs, set, score
            \item serveur, gagnant
            \item nombre de coups
            \item service : zone + lateralite
            \item sequence d'effets / lateralites
            \item dernier coup, derniere zone, resultat
        \end{itemize}
    \end{block}
    
    \begin{alertblock}{Idee cle}
        Un point n'est plus seulement une ligne de CSV.
        Il devient une \textbf{representation vectorielle du sens du point}.
    \end{alertblock}
    }
\end{frame}

\begin{frame}{Recherche semantique : index et calcul vectoriel}
    {\small
    \begin{block}{Document Elasticsearch}
        \begin{itemize}
            \item \textbf{metadonnees} : \texttt{match\_id}, \texttt{winner}, \texttt{serveur}, \texttt{set\_num}...
            \item \textbf{description} textuelle
            \item \textbf{embedding} dense 384D
        \end{itemize}
    \end{block}

    \begin{block}{Calcul cote requete}
        \begin{itemize}
            \item texte utilisateur $\rightarrow$ embedding de requete
            \item meme espace vectoriel que les points indexes
            \item comparaison requete / points par \textbf{proximite vectorielle}
        \end{itemize}
    \end{block}
    
    \begin{block}{kNN dans Elasticsearch}
        \begin{itemize}
            \item champ utilise : \texttt{embedding}
            \item \texttt{k = nombre de resultats}
            \item \texttt{num\_candidates = 10 * k}
            \item objectif : retrouver les \textbf{k plus proches voisins}
        \end{itemize}
    \end{block}

    \begin{alertblock}{Interpretation}
        Deux points proches dans l'espace vectoriel
        $\Rightarrow$ descriptions proches
        $\Rightarrow$ \textbf{sens proche}, meme sans mots identiques.
    \end{alertblock}
    }
\end{frame}

\begin{frame}{Recherche semantique : recherche hybride}
    {\small
    \begin{block}{Combinaison de deux signaux}
        \begin{itemize}
            \item \textbf{signal vectoriel} : similarite semantique
            \item \textbf{signal textuel} : \texttt{multi-match} sur
            \texttt{description}, \texttt{sequence\_effets}, \texttt{sequence\_coups}
            \item \textbf{poids vectoriel = 0.7}
        \end{itemize}
    \end{block}

    \begin{block}{Filtres exacts conserves}
        \begin{itemize}
            \item \texttt{match\_id}
            \item \texttt{winner}
            \item \texttt{serveur}
            \item \texttt{set\_num}
            \item \texttt{faute\_type}, \texttt{dernier\_coup}, \texttt{service\_zone}...
        \end{itemize}
    \end{block}

    \begin{alertblock}{Interet}
        \begin{itemize}
            \item moins sensible aux mots exacts
            \item robuste aux reformulations
            \item compatible avec les filtres metier
            \item meilleur compromis : \textbf{sens + contraintes}
        \end{itemize}
    \end{alertblock}
    }
\end{frame}

\end{document}
